\documentclass[11pt]{article}
\usepackage[margin=1in]{geometry}
\usepackage{times}
\usepackage{hyperref}
\hypersetup{colorlinks=true, linkcolor=black, urlcolor=blue, citecolor=black}
\title{The Center of Mass Does Not Negotiate: The Dean Drive and the Recurring Appeal of Oscillating-Mass Propulsion}
\author{Flyxion \\ \small Independent Researcher}
\date{}
\begin{document}
\maketitle

\begin{abstract}
The Dean Drive, patented by Norman Dean in the 1950s and periodically reinvented under other names since, proposes that an arrangement of eccentrically rotating or vibrating masses within a closed system can produce net unidirectional thrust without expelling any propellant. This paper argues that the device belongs to a broad and recurring family of internal-mass propulsion claims whose common flaw is not a subtle error in a specific gearing scheme but a direct violation of the conservation of linear momentum for an isolated system, a symmetry with no exception carved out for asymmetric timing, mechanical rectification, or any of the variations these devices have proposed across seven decades of reinvention.
\end{abstract}

\section{Introduction}

Norman Dean's device, patented in 1959 and 1961, used a system of rotating eccentric weights coupled through a mechanism intended to apply a stronger force during one phase of the rotation than the opposite phase, producing, in Dean's account, a net thrust in a fixed direction relative to the device housing, with no reaction mass expelled. Dean claimed demonstrations of the device lifting its own weight or reducing effective weight during operation, though these demonstrations were never conducted under independently controlled conditions, and no working replica has produced net thrust under laboratory scrutiny since. Numerous descendants of the same basic idea, internal masses moved in some cyclic, asymmetric pattern intended to produce net motion of the whole system, have appeared under different names in the decades since, including various self-described inertial propulsion and reactionless drive patents.

\section{Why Internal Asymmetry Cannot Do What Is Claimed}

For a closed system with no external forces acting on it, the center of mass of the system cannot accelerate, regardless of how the masses within the system move relative to each other. This is a direct consequence of Newton's third law applied to the internal forces between the system's components: every internal force has an equal and opposite reaction internal to the same system, so the vector sum of all internal forces is identically zero, and the system's total momentum, and therefore the motion of its center of mass, cannot change. This holds regardless of how asymmetric the timing of the internal motion is, how eccentric the rotating masses are, or how cleverly the mechanical linkage is arranged to make one phase of the cycle look different from another to a casual observer. A device that appears to move forward during one part of its cycle and only partially back during the other is, without exception, moving its housing relative to an internal counterweight in a way that leaves the system's overall center of mass exactly where conservation of momentum requires it to remain, with any apparent net displacement attributable to the device's mounting, friction with the supporting surface, or measurement over an incomplete cycle.

\section{Why Friction Is the Quiet Explanation Every Time}

The reason variations on this idea recur rather than being permanently settled by the momentum argument is that a device resting on a surface with static and kinetic friction is not, technically, a closed system isolated from external forces, since the supporting surface can exert horizontal frictional forces during part of the device's cycle and not during another, if the internal masses' acceleration profile changes the normal force against the surface asymmetrically across the cycle. This can produce genuine net creeping motion across a table, and several historical "successful" demonstrations of Dean-Drive-type devices are now understood to have relied on exactly this effect, disappearing entirely when the same device is tested on a sufficiently low-friction surface, suspended, or operated in free fall, conditions under which every rigorously tested version of this device family has failed to produce any net motion.

\section{The Entropic Gravity Framing}

This family of claims does not typically invoke gravitation directly, proposing instead a violation of momentum conservation through mechanical means, but the versions that do claim an accompanying weight reduction effect run into the same absence of coupling mechanism addressed elsewhere in this series: an oscillating internal mass, however configured, has no established or theoretically motivated channel into the large-scale ordered-configuration redistribution that produces gravitational curvature, on the entropic account or on the standard general-relativistic one, and the weight-reduction claims attached to some Dean Drive descendants add an unsupported gravitational layer on top of an already-refuted momentum claim rather than offering independent support for it.

\section{Objections}

Proponents sometimes argue that conservation of momentum, while valid in the timescale-averaged sense, might permit a net displacement over an incomplete or specifically truncated operating cycle, since the theorem strictly applies over a full period of the internal motion. This is occasionally true in a narrow technical sense for systems interacting with an external friction or drag force during only part of their cycle, which is precisely the friction-mediated creeping effect described above, and is not the reactionless, friction-independent thrust the original claims describe; no version of this device family has demonstrated net motion once external asymmetric friction and mounting effects are controlled for.

\section{Conclusion}

The Dean Drive and its many descendants propose a mechanical route around a conservation law with no carved-out exception for internal timing asymmetry, and the demonstrations that have appeared to succeed have, on closer examination, relied on exactly the external friction interaction that the underlying physics predicts would be necessary to produce real motion, which is to say they have worked, when they have worked at all, for reasons entirely outside the mechanism their inventors described.

\end{document}
